% titel: Blatt 2, A5: Verständnisfrage -- Newton'sche Gesetze
% skript: Kap. Dynamik, Newton'sche Gesetze

% Drei Axiome, Schwerelosigkeit, Actio = Reactio.

Die Newton'sche Mechanik beruht auf drei Axiomen, die die Wirkung von
Kräften auf Bewegungen beschreiben.

\begin{werkzeugkasten}[Werkzeuge]
\textbf{1. Axiom (Trägheit):} Ohne resultierende Kraft bleibt ein Körper
in Ruhe oder bewegt sich gleichförmig:
\begin{equation}
  \vec F_{\mathrm{ges}}=0\;\Leftrightarrow\;\vec a=0.
\end{equation}

\textbf{2. Axiom (Grundgleichung):} Die Beschleunigung ist proportional zur
Kraft und umgekehrt zur Masse:
\begin{equation}
  \vec F=m\vec a.
\end{equation}

\textbf{3. Axiom (Actio = Reactio):} Zwei Körper wirken mit gleich grossen,
entgegengesetzten Kräften aufeinander:
\begin{equation}
  \vec F_{12}=-\vec F_{21}.
\end{equation}
Die Kräfte wirken auf \emph{verschiedene} Körper.
\end{werkzeugkasten}

\begin{werkzeugkasten}[Inertialsysteme und Schwerelosigkeit]
Ein \textbf{Inertialsystem} ist ein System, in dem das 1. Axiom gilt
(unbeschleunigt). Die ISS und die Astronauten in ihr fallen beide frei
um die Erde; das System ist kein Inertialsystem (beschleunigt auf der
Kreisbahn), aber die \emph{Erscheinung} der Schwerelosigkeit entsteht,
weil alle Teile gemeinsam gleiche Bahnbeschleunigung haben und keine
Normalkraft wirkt.
\end{werkzeugkasten}

\begin{werkzeugkasten}[Vorgehen]
\begin{itemize}
  \item Für jeden Körper einzeln die Kräfte eintragen (Freikörperbild).
  \item ``Schwerelos'' heisst: es wirkt \emph{dennoch} die Gewichtskraft;
        nur die scheinbare (Normal-)Kraft fehlt.
  \item Beim Pferd-Wagen-Problem: Actio=Reactio wirkt auf verschiedene
        Körper, daher heben sich die beiden Kräfte nicht auf.
\end{itemize}
\end{werkzeugkasten}

\begin{achtung}[Typische Fallen]
\begin{itemize}
  \item Actio und Reactio \emph{nicht} addieren: sie betreffen verschiedene
        Körper, das 2. Axiom gilt für den einzelnen Körper.
  \item ``Keine Kraft'' $\neq$ ``kraftlos''; die ISS bewegt sich nicht
        gleichförmig (Kreisbahn).
\end{itemize}
\end{achtung}
